\scanpage{441}
\mspage{223}

of topological automorphisms of $\widetilde U$ compatible with the
equivalence relation \unsure{defined} by $\widetilde U\unsure{\longrightarrow|-}U$, and
inducing on $U$ the \unsure{automorphism} $\mathrm{id}_U$ or $\sigma$.

\deleted{For} Then

a) (For the record) $U^\sigma$ is a closed submanifold
of $U$ of dimension 1.

b) For $\forall\ x\in U^\sigma$, one finds a
\inserted{$\pi$-}conjugacy class of \deleted{elements \unsure{of}}\inserted{splittings} of the extension
$\mathcal E$, in the usual way, whence a
map
\deleted{(\unsure{4|5})}
\[
 \begin{aligned}
 U^\sigma\longrightarrow\mathrm{Sc}(\mathcal E,\mathbb Z/2)
 &=\text{set of classes under}\\
 &\quad\text{$\pi$-conjugation of splittings of $\mathcal E$}\\
 &\quad\text{\unsure{over} $\mathbb Z/2$, or again of}\\
 &\quad\text{elements of order $2$ of}\\
 &\quad\text{$\mathcal E$ which are not in $\pi$.}
 \end{aligned}
 \tag{5}
\]

This map factors through a
\uline{bijection}
\[
 \pi_0(U^\sigma)\simeq\mathrm{Sc}(\mathcal E,\mathbb Z/2)
 \tag{6}
\]

c) Let $Z_i$ be a connected component of $U^\sigma$,
corresponding to a \deleted{splitting} $\sigma_i\in\mathcal E^-$ \deleted{such}\inserted{of}
order $2$. Then \deleted{\illegible{1} $\pi$\illegible{1}}
\[
 \pi^{\sigma_i}=\{1\}\quad\text{\unsure{sss|iff}}\quad Z_i\simeq\mathbb R,
 \quad\deleted{\unsure{\text{and}}}
 \tag{7}
\]

d) If $Z_i\not\simeq\mathbb R$ i.e. \deleted{\illegible{4} $Z_i$} $Z_i\simeq\unsure{S^1|\mathbb S^1}$, then,
for a chosen orientation of $Z_i$, denot-

\scanpage{442}
ing by $g_i$ the element of $\pi=\pi_1(U)$
that it defines (defined up to conjugation), and
by $\varphi_i:\mathbb Z\longrightarrow\pi$ the homomorphism $\varphi_i(n)=g_i^n$ from
$\mathbb Z$ to $\pi$, one has:

$1^\circ)$ $\varphi_i$ is injective

$2^\circ)$ after possibly replacing $\varphi_i$ by a conjugate,
one has
\[
\unsure{(9)|(8)}\qquad \pi^{\sigma_i}=\operatorname{Im}\varphi_i
\]

\uline{Proof} One finds the map (5) by
\deleted{\unsure{taking}}, for $\forall\ x\in U^\sigma$, the \inserted{\unsure{canonical}} operation of
$\mathbb Z/2$ on the universal covering $\widetilde U(\unsure{x})$ pointed
at $x$, and by \deleted{choosing} taking the
\unsure{corresponding} operations on $\widetilde U$, \unsure{deduced} via
the isomorphisms $\widetilde U\simeq\widetilde U(x)$. One sees that
the \deleted{\unsure{only}} operations $\sigma'$\deleted{\illegible{1}} obtained on $\widetilde U$ \inserted{(for fixed $x$)}, are
those for which \deleted{\illegible{1}} $\widetilde U^{\sigma'}$ has a
point over $x$ --- hence the set
of $x$ which give rise to \deleted{\unsure{a}}\inserted{the}
class of a $\sigma'$, are the elements of
the image of $\widetilde U^{\sigma'}$ under the projection
$\widetilde U\longrightarrow U$. Since by the lemma $\widetilde U^{\sigma'}$
is nonempty and connected, it follows that

\scanpage{443}
\mspage{224}
its image in $U^\sigma$ is so too --- we are going to
see that its image is exactly a connected
component of $U^\sigma$ --- which \unsure{will prove} both
the factorizability of (5) through \deleted{\unsure{the}} $\pi_0(U^\sigma)$
(\unsure{rather} trivial in any case) and the fact that
the resulting map (6) is \uline{bijective}.

Let $\widetilde Z_i$ be a universal \unsure{cover\textsuperscript{ing}|covering} of $Z_i$,
and let us take the \unsure{corresponding} universal covering
$\widetilde U_i$ of $U$, so that one has
\margin{\deleted{\diagram{\begin{tikzcd}[ampersand replacement=\&]
\widetilde Z_i \arrow[r] \arrow[d] \& \widetilde U_i \arrow[d] \\
Z_i \arrow[r] \& U
\end{tikzcd}}}}
\unsure{(8)|(9)}
\diagram{\begin{tikzcd}[ampersand replacement=\&]
\widetilde Z_i \arrow[r] \arrow[d] \& \widetilde U_i \arrow[d] \\
Z_i \arrow[r] \& U
\end{tikzcd}}
and by functoriality $\sigma$ also operates on this
diagram (operating \unsure{trivial\textsuperscript{ly}} on $Z_i,\widetilde Z_i$),
let $\sigma_i$ be its operation on $\widetilde U_i$. One
thus sees that $\widetilde Z_i$ maps into $\widetilde U_i^{\sigma_i}$,
which therefore maps \uline{onto} $Z_i$ --- which already finishes
proving b).

\deleted{Let us suppose now that $Z_i$ is \unsure{simpl\textsuperscript{y}|simply} connected
i.e. $Z_i\simeq\mathbb R$, $\widetilde Z_i=Z_i$ --- thus $\widetilde Z_i\longrightarrow\widetilde U_i$ is
\unsure{a} section of $\widetilde U_i$.}
\inserted{Let us consider \deleted{\unsure{the image}} \inserted{the image} \deleted{\unsure{induced}} of $\widetilde Z_i$ in
$\unsure{\widetilde U_i|\widetilde U}\mid Z_i$, this is an \uline{open} and closed subset}
(as the image of a morphism of étale \unsure{cover\textsuperscript{ings}|coverings}
over the same component $Z_i$), so it is

\scanpage{444}
\unsure{fortiori} an open and closed subset of
$\widetilde U_i^{\sigma_i}$, and since this space is connected,
it is equal to it.
\deleted{When}
\deleted{\[
(10)\qquad \widetilde U_i^{\sigma_i}=\operatorname{Im}(\widetilde Z_i\longrightarrow\widetilde U_i)
\]}
\inserted{Moreover $\widetilde Z_i\longrightarrow\widetilde U^{\sigma_i}$ makes \deleted{\illegible{1}}\inserted{of}
$\widetilde Z_i$ an étale covering of
its image $\widetilde U_i^{\sigma_i}$ in $\widetilde U_i\mid Z_i$, $\ast$
and since $\widetilde U_i^{\sigma_i}$ is
\unsure{simpl\textsuperscript{y}|simply} connected,
one finds \unsure{final\textsuperscript{ly}|finally}
\[
(10)\qquad \widetilde Z_i\overset{\sim}{\longrightarrow}\widetilde U_i^{\sigma_i}
\]}

When $Z_i$ is \unsure{simpl\textsuperscript{y}|simply} connected i.e.
$\widetilde Z_i\simeq Z_i$, this also means that
$\widetilde U_i^{\sigma_i}\longrightarrow Z_i$ is a homeomorphism,
(thus $\widetilde Z_i=Z_i\longrightarrow\widetilde U_i^{\sigma_i}$ is the inverse
\unsure{homeomorphism}). Since, for $x\in Z_i$,
\deleted{\unsure{the elements}} and $\widetilde x\in(\widetilde U_i^{\sigma_i})_x$, the
$\widetilde x'$ of $\widetilde U_i^{\sigma_i}$ over $x$ are \deleted{\illegible{1}} \inserted{\unsure{the elements}} of
the form $\widetilde x\cdot\gamma$, with $\gamma\in\pi_i^{\sigma_i}$
($\pi_i=\operatorname{Aut}(\widetilde U_i/U)$) it follows that one
indeed has
\[
(11)\qquad \pi_i^{\sigma_i}=1
\]
which is essentially formula (7).

In the case \deleted{general, \unsure{denoting}} \inserted{$Z_i$ not \unsure{simpl\textsuperscript{y}|simply} connected,} setting
\[
(12)\qquad T_i=\operatorname{Aut}(\widetilde Z_i/Z_i)\underset{\unsure{f}}{\simeq}\pi_1(Z_i;\widetilde Z_i)
\]
one finds a \unsure{canonical} homomorphism
\[
(13)\qquad \varphi_i:T_i\longrightarrow\pi_i
\]
and injectivity in (10) is equivalent to
(13) being injective. Of course,

\scanpage{445}
\mspage{225}
the choice of a generator $g_i^{\circ}$ of $T_i$ amounts to
the choice of an orientation of $Z_i$, and
$g_i=\varphi_i(g_i^{\circ})\in\pi_i$ is then the corresponding element
of $\pi_1(U)$ in question in
the statement (less precise, \unsure{since} it only speaks
of a conjugacy class). It is clear
that $\sigma_i$ (\unsure{operating} \unsure{trivial\textsuperscript{ly}} on $T_i$, and on $\pi_i$
by transport of structure) commutes with
the injective homomorphism $\varphi_i$, hence \deleted{$\varphi_i(T_i)\subset$} $\varphi_i$
induces \deleted{$\varphi_i$} $T_i\longrightarrow\pi_i^{\sigma_i}$. I say that
this is in fact an \uline{isomorphism}
\[
\text{(14)}\qquad \varphi_i:T_i\xrightarrow{\sim}\pi_i^{\sigma_i}
\]
-- it remains to prove surjectivity. But
if $\gamma\in\pi_i^{\sigma_i}$, then $\widetilde{x}\gamma\in\widetilde{U}_i^{\sigma_i}$, hence
$\widetilde{x}\gamma\in\varphi_i(\widetilde{Z}_i)$ which means that
$\widetilde{x}\gamma=\widetilde{x}\varphi_i(g)$ \deleted{\unsure{hence}} for $g\in T_i$ hence
$\gamma\in\operatorname{Im}\varphi_i$, q.e.d.

\uline{Corollary} 1 Let $\sigma_i\in\mathcal{T}_g$ $(0\leq i\leq g)$,
corresponding to an anti-involution $\sigma_i$ of
$X_g$ such that $\operatorname{card}\pi_0(X_g^{\sigma_i})=i$. Then
\margin{and $X_g/\sigma$ nonorientable....}
a) There are exactly $i$ \inserted{$\Pi_g$}-conju-

\scanpage{446}
gacy classes of actual automorphisms of order $2$
of $\Pi_g$, in the class $\sigma_i$ (which proves
that if $i\neq j$, $\sigma_i$ and $\sigma_j$ are not
conjugate in $\mathcal{T}_g$ \deleted{\illegible{1}} \dots)

b) If $u_i\in\mathcal{G}_g=\operatorname{Aut}_{\mathrm{lac}}(\Pi_g)$ is \inserted{of order $2$,} in the
class $\sigma_i$, then \deleted{$\unsure{\Pi_g}$}
\[
\text{(14)}\qquad \Pi_g^{\unsure{u_i}}\simeq\mathbb{Z}
\]

c) If $u_i,u'_i\in\mathcal{G}_g$, \deleted{\unsure{satisfy} \unsure{both} anti-}
\inserted{are of order $2$}
\deleted{involutions} of class $\sigma_i$, then $\exists\,\unsure{h}\in\mathcal{G}_g^+$
such that
\[
\text{(15)}\qquad u'_i=hu_ih^{-1}=\operatorname{int}(h)u_i
\]
(\uline{N.B.} \unsure{Necessarily}, the image of $h$ in $\mathcal{T}_g^+$
will be in $(\mathcal{T}_g^+)^{\sigma_i}=\operatorname{Centr}_{\mathcal{T}_g}(\sigma_i)^+$.)

\uline{Proof} a) and b) are special
cases of the theorem, applied to
an anti-involution $\sigma_i$ of $X_g$, with
$\pi_0(X_g^{\sigma_i})$ of cardinality $i$. Let $Y_{\unsure{g,i}}=X_g^{\sigma_i}$,
a compact surface \inserted{with boundary} nonorientable \deleted{of} connected
of type $(g-i,i)$, it is well known,
and immediate, that the group of auto-

\scanpage{447}
\mspage{226}
morphisms \unsure{of the|of}\deleted{$\unsure{Y}$} such a manifold is transitive
on $\pi_0(\partial Y)$, hence on \unsure{lifting} to $(X_g,\sigma_i)$, that
the group of automorphisms of $(X_g,\sigma_i)$ is transitive
on $\pi_0(X_g^{\sigma_i})$, \deleted{\illegible{2}} which one can interpret
also as \unsure{the set} of $\Pi_g$-conjugacy classes
of $u_i$ as in b), c). This
therefore implies that there exists $\dot{h}\in\unsure{\mathcal{T}_g^+|\mathcal{G}_g^+}$, \inserted{commuting with $\sigma_i$,} such that
-- denoting by $h$ a \unsure{lifting|lifting} in $\mathcal{G}_g^+$ --
one has $u'_i$ $\unsure{\Pi|\pi}$-conjugate to $\operatorname{int}(h).u_i$, which
also means that one can after possibly modifying
$h$), choose it so that one has (15)

\uline{Corollary} 2 Under the conditions of the preceding
corollary, \unsure{setting} $\Pi_g^{u_i}=T$ $(\simeq\mathbb{Z})$, one
has a diagram of exact sequences
\[
\text{(16)}\qquad
\diagram{\unsure{
\begin{tikzcd}[ampersand replacement=\&,column sep=large]
1 \arrow[r] \& T \arrow[r] \arrow[d,equal] \&
\deleted{\unsure{\operatorname{Centr}_{\mathcal{G}_g}}}\mathcal{G}_g^{+u_i}
\arrow[r] \arrow[d] \&
\mathcal{T}_g^{+\sigma_i}
\arrow[d,hook] \quad \deleted{\unsure{\longrightarrow 1}}
\\
1 \arrow[r] \& T \arrow[r] \&
\mathcal{G}_g^{\deleted{+}u_i} \arrow[r] \&
\mathcal{T}_g^{\deleted{+}\sigma_i}
\quad \deleted{\unsure{\longrightarrow 1}}
\end{tikzcd}
}}
\]
and the index of $\mathcal{G}_g^{+u_i}$ in $\mathcal{T}_g^{+\sigma_i}$
(and of $\mathcal{G}_g^{u_i}$ in $\mathcal{T}_g^{\sigma_i}$) is $i$.
Since the last two vertical arrows
are inclusions of subgroups of index $2$,
it suffices to treat the\deleted{\unsure{es}} assertion concerning
$\mathcal{G}^+$, $\mathcal{T}^+$, \deleted{which \illegible{2} \inserted{\illegible{1}} \illegible{1}}

\scanpage{448}
\deleted{1°) $\mathcal{G}_g^{+u_i}\longrightarrow\mathcal{T}_g^{+\sigma_i}$ is surjective. Let
$\unsure{\dot\gamma}\in\mathcal{T}_g^{+\sigma_i}$, then $\unsure{u_i'}=\operatorname{int}(\gamma).u_i$ satisfies the
conditions of c), hence $\exists h\in\mathcal{G}_g^{+\deleted{u_i}}$ such that
\deleted{\illegible{1}} $u_i'=\operatorname{int}(h)u_i$ hence $\operatorname{int}(\unsure{h\gamma})u_i=u_i$ i.e.
$\unsure{h\gamma}\in$.}

Now $\mathcal{T}_g^{+\sigma_i}$ operates \inserted{transitively} on the set of
$\Pi$-conjugacy classes of the $u_i'$, by
c), \unsure{on the other hand} the stabilizer \inserted{in $\mathcal{T}_g^{+\sigma_i}$}, for this
action, of $u_i$, consists of the $\dot\gamma\in\mathcal{T}_g^{+\sigma_i}$
such that \deleted{\unsure{there exists}} $\operatorname{int}(\gamma)u_i$ is $\Pi$-conjugate
to $u_i$, i.e. such that there exists $\alpha\in\Pi$, with
$\operatorname{int}(\gamma)u_i=\operatorname{int}(\alpha)u_i$ i.e. $\alpha^{-1}\gamma\in\mathcal{G}_g^{+u_i}$, which
means that $\dot\gamma$ is in the image of $\mathcal{G}_g^{+u_i}$,
q.e.d.

\margin{\deleted{\unsure{as}}}
\deleted{\uline{Corollary 3} Let $U$ be a \deleted{compact} surface
\inserted{\deleted{\illegible{1}} \unsure{0-connected}, not \unsure{homeomorphic} to the real projective plane.}
with boundary, $\Pi=\pi_1(U)$, (For \deleted{\illegible{1}} $\Pi$
to have torsion elements, it \unsure{is necessary and sufficient}
that $U$ be nonorientable and \deleted{$\partial U\ne\varnothing$. Then}
\inserted{that $\partial U$ have at \unsure{least} one compact connected component.}
\deleted{\illegible{1}} The torsion elements of
\deleted{\unsure{when} $U$ \illegible{2} \unsure{orientable},}
$\Pi$ \deleted{\unsure{as}} are of order 2, and their conjugacy
classes are in 1-1 correspondence with the
elements of $\pi_0(\partial U)$ which \unsure{correspond} to compact \unsure{connected} components.}

\scanpage{449}
\mspage{227}
I feel like constructing a \inserted{geometric} figure where
one can bring out simultaneously
topological anti-involutions
which \deleted{are} give rise to the
$\sigma_i\in\unsure{\mathcal{T}_g^{\natural}|\mathcal{T}_g^{\sharp}}$ (up to conjugation\unsure{..}), and to the various
$u_i$ \inserted{$\in\mathcal{G}_g$} associated with a $\sigma_i$ (which will then be
easy, by the general recipe). We
know that $\sigma_i$ \inserted{$\in\mathcal{T}_g$} is obtained by regarding
an $X_g$ as the oriented double
of a $Y=Y_{g-\deleted{\unsure{2}}i,i}$, \unsure{thus starting}
from the \deleted{\illegible{1}} real projective plane $Y_{\deleted{\unsure{g}}0}$, by
\margin{\deleted{\unsure{$D_0$}}}
making $g$ holes in it, \deleted{of which} one \inserted{\deleted{\illegible{1}}} fills
$\unsure{g-i}$ back in with Möbius bands,
leaving the $i$ other holes as they are.
Let $D_j$ ($1\le j\le g$) be the disjoint
closed disks \unsure{corresponding} to the ``holes, thus
\deleted{\unsure{Then}}
\[
(17)\qquad \boxed{V_{0,j}=Y_0\setminus\bigcup_j \overset{\circ}{D}_j}
\]
is a manifold with boundary
(nonorientable), contained both in
$Y_0$, and in $Y=Y_{g-i,i}$, and even
coinciding with $Y$ \unsure{at the points} of
\[
\boxed{\partial Y=\bigcup_{\unsure{g-i+1\le j\le g}}\partial D_j}.
\]
\deleted{Thus} Let $X_0$ be the oriented sphere which \unsure{covers} $Y_0$,

\scanpage{450}
and $U_0$ \deleted{\illegible{1}} $=X_0\vert V_0$, its restriction over $\bigcup D_j$,
it is thus \inserted{the complement of} a union of $2g$ disks
\[
(18)\quad \deleted{\unsure{U_0}}\quad U_0\deleted{\unsure{U_g^i}}=X_0\setminus\bigcup_{1\le j\le g}\deleted{D}\Delta_j,
\]
where $\Delta_j\longrightarrow D_j$ is a trivial $2\deleted{\unsure{g}}$-sheeted
covering of $D_j$ ($\Delta_j=\Delta'_j\amalg\Delta''_j\simeq D_j\times\unsure{\varepsilon_j|\mathcal{E}_j}$,
where $\unsure{\varepsilon_j|\mathcal{E}_j}$ is a set with 2 elements, which is identified
with the set of the two orientations of
$D_j$\unsure{...}). \deleted{One \unsure{constructs} \illegible{1} so that}
\inserted{Let \unsure{on the other hand}}
\deleted{\illegible{1}} $M_j$ \inserted{($1\le j\le g-i$)} be the Möbius band whose boundary \unsure{has|.}
been glued back to $V_0$ by a homeomorphism
\[
(19)\qquad \partial M_j\simeq\partial D_j
\]
\unsure{let|over}
\[
(20)\qquad M=\coprod_{1\le j\le g-i} M_i
\]
so that
\[
(21)\qquad Y_{g\unsure{-|+}i,i}=V_0\amalg_{\partial M}M.
\]

\deleted{One sees that one has a continuous map
\[
U_0\longrightarrow\unsure{X_{g-i,i}|X_{g,i}}
\]
which induces a homeomorphism
\[
\operatorname{Int}(U_0)\longrightarrow Y_{g-i,i}\setminus(M\cup\partial Y)
\]
More precisely, the inverse image}

Let then $X=X_g$ be the oriented \unsure{double|double}
of $Y_{g-i,i}$ \unsure{defined} by (21), let $\widetilde M=\coprod_{1\le j\le g-i}\widetilde M_j$
be the inverse image of $M$ in it, \deleted{\unsure{$X_g$} \illegible{4}}
\inserted{an étale covering of degree 2 of $M$,} \deleted{\illegible{3}}
\[
(\unsure{22|21})\qquad \deleted{X_g\unsure{\illegible{1}}}\quad X_{g,i}=X_g\setminus\operatorname{Int}(\widetilde M),
\]
